% !TEX program = lualatex
% ECON 201, Worksheet 13: Strategic Price Setting.
% Compile with LuaLaTeX so the PDF is tagged. Build from the course root:
%   latexmk -lualatex -cd worksheets/worksheet13.tex && latexmk -c -cd worksheets/worksheet13.tex
% Every number here is computed and printed by slides/figures/make_lecture13.py
% (Wanda and Kit with 14 loyal customers).
\DocumentMetadata{
  pdfstandard=UA-2,
  pdfversion=2.0,
  lang=en-US,
  tagging=on
}
\documentclass[11pt]{article}

\usepackage[letterpaper, margin=0.75in, top=0.5in, bottom=0.45in]{geometry}
\usepackage{fontspec}
\setmainfont{Fira Sans}
\usepackage{amsmath}
\usepackage{amssymb}
\usepackage{graphicx}
\usepackage{booktabs}
\usepackage{array}
\usepackage{xcolor}
\usepackage{needspace}
\usepackage[most]{tcolorbox}
\definecolor{cream}{HTML}{FDF3EA}
\definecolor{accent}{HTML}{BF5700}
\usepackage{titlesec}
\titleformat{\section}{\large\bfseries\color{accent}}{}{0pt}{}
\titlespacing*{\section}{0pt}{4pt}{2pt}
\usepackage{hyperref}
\hypersetup{pdftitle={ECON 201 Worksheet: Strategic Price Setting}, pdfauthor={Div Bhagia}, hidelinks}
\pagestyle{empty}
\setlength{\parindent}{0pt}
\setlength{\parskip}{2pt}

% Ruled answer space: n lines. Plain TeX loop; pgffor's \foreach breaks the
% enumerate counter when used inside an item.
\newcount\answerlinecount
\newcommand{\answerlines}[1]{%
  \par\vspace{1pt}\answerlinecount=#1
  \loop\ifnum\answerlinecount>0
    \rule{\linewidth}{0.3pt}\par\vspace{5pt}%
    \advance\answerlinecount by -1
  \repeat
  \vspace{-5pt}}

% Ruled grids rather than booktabs tables: students write in the cells, so
% every cell gets visible borders and the blank cells get extra height.
\newcolumntype{C}{>{\centering\arraybackslash}p{2.6cm}}
\newcommand{\blankrow}{\rule{0pt}{1.5em}}

\begin{document}

{\LARGE\bfseries\color{accent} Worksheet: Strategic Price Setting}\hfill{\large ECON 201}

% Terms sit on the cream panel, where the brand orange falls below 4.5:1;
% the darker orange keeps WCAG AA contrast (see CLAUDE.md).
\definecolor{accentdark}{HTML}{8F4100}
\newcommand{\term}[1]{\textcolor{accentdark}{\textbf{#1}}}

% Definitions panel list: flush left, compact (see worksheet09).
\newenvironment{deflist}
  {\begin{list}{\textbullet}{%
     \setlength{\leftmargin}{1.1em}\setlength{\labelwidth}{0.7em}%
     \setlength{\labelsep}{0.4em}\setlength{\itemindent}{0pt}%
     \setlength{\listparindent}{0pt}\setlength{\topsep}{2pt}%
     \setlength{\partopsep}{0pt}\setlength{\parsep}{0pt}%
     \setlength{\itemsep}{2pt}}}
  {\end{list}}

\begin{tcolorbox}[colback=cream, colframe=accent, boxrule=0.6pt, arc=2pt,
  left=6pt, right=6pt, top=4pt, bottom=4pt, title=Definitions,
  fonttitle=\bfseries, coltitle=white]
\begin{deflist}
\item \term{Best response}: the strategy that gives a player the highest
  pay-off, given what the others choose.
\item \term{Nash equilibrium}: a set of strategies, one for each player,
  where each player's strategy is a best response to the others'.
\item \term{Dominant strategy}: a strategy that is a player's best response
  whatever the others choose.
\item \term{Prisoners' dilemma}: a game where each player has a dominant
  strategy, but another outcome gives every player a higher pay-off.
\end{deflist}
\end{tcolorbox}

\needspace{5\baselineskip}\section{Wanda and Kit, with fewer loyal customers}

Each rental costs \$10, and each firm picks $H = \$36$ or $L = \$20$. Now each firm has \textbf{7 loyal} customers, \textbf{13 switchers}, and \textbf{10 bargain hunters}. When one firm is cheaper, the other firm's switchers and bargain hunters come to it.

\vspace{2pt}
1. Fill in the last two columns of Wanda's customers and profit. The first two are done for you.

\begin{center}
\renewcommand{\arraystretch}{1.25}
\begin{tabular}{|l|c|c|C|C|}
\hline
\textbf{Wanda's price, Kit's price} & $H, H$ & $L, L$ & $L, H$ & $H, L$ \\
\hline
Her loyal customers   & 7  & 7  & \blankrow & \blankrow \\ \hline
Her switchers         & 13 & 13 & \blankrow & \blankrow \\ \hline
Kit's switchers       & 0  & 0  & \blankrow & \blankrow \\ \hline
Her bargain hunters   & 0  & 10 & \blankrow & \blankrow \\ \hline
Kit's bargain hunters & 0  & 0  & \blankrow & \blankrow \\ \hline
\textbf{Total customers} & 20 & 30 & \blankrow & \blankrow \\ \hline
\textbf{Profit}, $(P - 10) \times$ total & 520 & 300 & \blankrow & \blankrow \\ \hline
\end{tabular}
\end{center}

2. Kit's numbers mirror Wanda's. Fill in the pay-off matrix, with Wanda's profit first in each cell.

\begin{center}
\renewcommand{\arraystretch}{1.3}
\begin{tabular}{|c|C|C|}
\hline
 & \textbf{Kit picks $H$} & \textbf{Kit picks $L$} \\
\hline
\textbf{Wanda picks $H$} & 520, 520 & \blankrow \\
\hline
\textbf{Wanda picks $L$} & \blankrow & 300, 300 \\
\hline
\end{tabular}
\end{center}

3. If Kit picks $H$, Wanda earns \underline{\hspace{1.4cm}} at $H$ and \underline{\hspace{1.4cm}} at $L$, so her best response is \underline{\hspace{1cm}}. \\[4pt]
If Kit picks $L$, Wanda earns \underline{\hspace{1.4cm}} at $H$ and \underline{\hspace{1.4cm}} at $L$, so her best response is \underline{\hspace{1cm}}.

4. Does Wanda have a dominant strategy? \underline{\hspace{2cm}} \quad The Nash equilibrium is \underline{\hspace{2.2cm}}, with profits of \underline{\hspace{1.6cm}} each.

5. Is there another outcome where \emph{both} firms would earn more? Which one? What kind of game is this?
\answerlines{2}

\end{document}
